\documentclass[12pt,a4paper]{article}
\usepackage[utf8]{inputenc}
\usepackage[T1]{fontenc}
\usepackage[brazil]{babel}
\usepackage{amsmath,amssymb,amsthm}
\usepackage{geometry}
\geometry{margin=2.5cm}

\newtheorem{theorem}{Teorema}[section]

\title{Soma dos primeiros $n$ números naturais}
\author{}
\date{}

\begin{document}
\maketitle

\begin{theorem}
Para todo inteiro $n \geq 1$,
\[
  1 + 2 + 3 + \cdots + n = \frac{n(n+1)}{2}.
\]
\end{theorem}

\begin{proof}
Faremos indução sobre $n$.

\textbf{Base} ($n=1$): o lado esquerdo é $1$ e o direito é
$\dfrac{1\cdot 2}{2} = 1$. A igualdade vale.

\textbf{Hipótese de indução:} suponhamos que, para certo $k \geq 1$,
\[
  1 + 2 + \cdots + k = \frac{k(k+1)}{2}.
\]

\textbf{Passo indutivo:} queremos mostrar a igualdade para $k+1$.
Somando $k+1$ a ambos os lados da hipótese:
\[
  1 + 2 + \cdots + k + (k+1)
  = \frac{k(k+1)}{2} + (k+1)
  = \frac{k(k+1) + 2(k+1)}{2}
  = \frac{(k+1)(k+2)}{2}.
\]
O resultado é exatamente a fórmula com $n = k+1$. Pelo Princípio da
Indução Finita, a igualdade vale para todo $n \geq 1$.
\end{proof}

\end{document}
